\centersection{Abstract}

\begin{spacing}{1.5}
  This work investigates the application of the MeanFlow one-step generative model to image denoising tasks. The primary objective of this summer research is to understand the fundamental principles of Flow Matching generative modeling and its improved version, the MeanFlow model, through theoretical derivation and experiments. Based on this understanding, we specifically adapt the MeanFlow model for image denoising by modifying its initial distribution: from simple distributions like Gaussian to noisy images, retaining the model’s one-step generation capacity while serving the denoising purpose. Experiments on MNIST dataset show that a single function evaluation (NFE = 1) significantly reduces the noise, increasing the Peak Signal-to-Noise Ratio (PSNR) from approximately 6 dB to 22.6 dB. A second evaluation (NFE = 2) further improves the PSNR to 24 dB, while additional sampling steps yield limited gains. This work demonstrates that the MeanFlow framework is capable of performing effective image denoising at minimal sampling cost, positioning it as a practical alternative to conventional multi-step denoising methods. Although our work does not propose new algorithms or model frameworks, it explores the application of the MeanFlow framework to image denoising, verifying its feasibility for future image restoration tasks.
\end{spacing}

\clearpage
